\section{A sixth-degree tilt and cutoff seven}\label{sec:seven}

Retain \(r=e^{-x}\), \(p=n-1\), \(\rho=\sqrt{x^2+\pi^2}\), \(s=\pi/\rho\), and \(\zeta=x+\ii\pi\).
The main quantitative result of this section is the bound \(n\ge96\) with
the same five independent short-cycle weights in both rows.
For the cutoff-seven profile the normalized generating function and weighted
insertion formula from Section~\ref{sec:general} give
\begin{align*}
 B(z)&=\exp\!\left(\sum_{a=1}^{6}b_a z^a\right),
 & b_a&=(w_a-1)M_a/a!,\\
 Z&=\frac1n\sum_{a=1}^{6} C_a\rho^{a+1}Y_{n-a-1},
 & C_a&=(w_{a+1}-w_a)M_{a+1}/a!.
\end{align*}
Thus the shortest row in the insertion is \(n-7\). We prove
\(|Z|<|Y_n|\) at \(Y_p=0\) for all \(n\ge96\).

Write \(u_0=r/(1+r)^2\) and \(v_0=(1-r)/(1+r)\) as in Section~\ref{sec:six}.
The Eulerian recurrence gives \(A_6(t)=1+57t+302t^2+302t^3+57t^4+t^5\) and hence
\[
 M_7=-u_0v_0(1-60u_0+360u_0^2).
\]
This moment requires a new bound. At \(r=1/11\), \(|M_7|/r=258335/248832>1\).
In the variable \(v=v_0\in[0,1]\), let \(P(v)=17v/2-30v^3+45v^5/2\).
Then \(-M_7/u_0=P(v)\) and \(-M_7/r=(1+v)^2P(v)/4\).
For each of the 32 closed subintervals with endpoints \(h/32,(h+1)/32\),
apply the substitution \(v=(h+z)/32\). If the resulting power coefficients
are \(a_i\) in degree \(d\), the Bernstein coefficients are
\(\sum_{i=0}^{j}a_i\binom ji/\binom di\) for \(0\le j\le d\).
The Bernstein basis is nonnegative on \([0,1]\) and sums to one.
The maxima of the absolute coefficients are bounded by
\(152355739/83886080<11/6\) and \(1003440476713/962072674304<21/20\),
respectively. This proves \(|M_7|\le\min((11/6)u_0,(21/20)r)\), with the
endpoint \(r=0\) understood by continuity. All six tilt coefficients satisfy
\(|b_a|\le r/a!\) by the bounds through order six. For the drift,
\(|C_a|\le r/a!\) for \(a\le5\) and \(|C_6|\le(21/20)r/6!\).
Write \(\mathcal V(R)=\sum_{a=1}^{6}|b_a|R^a\) and
\(J=\sum_{a=1}^{6}|C_a|\rho^{a+1}/s\).

\paragraph{Compact interval \(0\le x\le16\).}
Use the adjustable two-pole bound \((*)\) from Section~\ref{sec:five},
at \(R=\lambda\rho\) with \(1<\lambda\le|x+3\ii\pi|/\rho\).
In each of the 128 closed eighth-grid intervals enclose \(r,u_0,v_0,\rho\)
by their endpoint bounds, using \(|M_6|\le u_0\) and both seventh-moment
bounds above. The rational certificate selects \(\lambda\) on the
grid of twentieths. It tests admissibility at the right endpoint,
where \((x^2+9\pi^2)/(x^2+\pi^2)\) is smallest. With \(D=|B(\zeta)|\),
the pole error divided by \(sD\) is at most
\[
 \epsilon=\max\!\left(\rho,\frac{\lambda\rho^2}{2\pi^2}\right)
 \frac{\rho}{\pi}\lambda^{-j}
 \exp\!\bigl(\mathcal V(\lambda\rho)+\mathcal V(\rho)\bigr).
\]
Use \(j\ge n-7\ge89\) and an upper bound for this expression on the cell.
The two-pole argument in Section~\ref{sec:six} yields
\(|Y_n|\ge(2-3\epsilon)sD\) at \(Y_p=0\) and
\(|Y_j|\le(2+\epsilon)D\) for the lagged rows.
Rational endpoint and exponential enclosures give a maximum of
\(\bigl(3\epsilon+(2+\epsilon)J/96\bigr)/2\) at most
\(246041/1250000<1\). The recorded upper bound allows every choice of
weights in the closed cube and proves the insertion inequality.

\paragraph{Middle interval \(16\le x\le3\sqrt p\).}
Let \(D_6(R)=\sum_{a=1}^6 R^a/a!\), \(K=D_6(16)/16^6=100117/47185920\)
and \(\alpha=51/50\). We have \(\rho\le\alpha x\),
\(\mathcal V(\rho)\le Kr\rho^6<1/200\), and
\(\exp(\mathcal V(2\rho))\le(1-64/200)^{-1}<2\).
Use \(157/50<\pi<22/7\), \(e^{-16}<1/8800000\), and the decrease
of \(x^m e^{-x}\) for \(x\ge16\) and \(m\le9\).
Split the tilt convolution at \(i=\lfloor j/2\rfloor\), for \(j\ge n-7\).
In the head, \(2(j-i+1)\ge n-5\), so the logistic bound of
Section~\ref{sec:six} gives \(2+2\rho^2/((n-5)\pi^2)\).
The total head error divided by \(s\) is therefore bounded by
\[
 \frac{K e^{-x}}{\pi(1-1/200)}
       \left(2\rho^7+\frac{2\rho^9}{(n-5)\pi^2}\right)
 <\frac1{16}.
\]
For the tail use \(|V_m|\le2\rho^2/x\) at all needed orders including zero.
Since \(2\rho^3/(\pi x)\le18\alpha^3p/\pi<7n\),
the doubled-radius bound makes the tail divided by \(s\) at most
\(14n\,2^{-\lfloor(n-7)/2\rfloor}<1/1000\), checked at \(n=96,97\)
and decreasing within each parity. Thus \(|Y_j-V_j|<s/4\).
We have \(ns^2\ge1\) since \(x^2\le9p<\pi^2p\).
At \(Y_p=0\) the bare gap \eqref{center} gives \(|V_n-cV_p|>s\),
where \(c=x/\rho\). The two errors give \(|Y_n|>s/2\), and all six
lagged rows satisfy \(|Y_j|<2+n/(n-6)+1/4<4\). Finally
\[
 J\le\frac{r}{\pi}\left(\sum_{a=1}^{5}\frac{\rho^{a+2}}{a!}
                        +\frac{21\,\rho^8}{20\cdot6!}\right)<2/5 .
\]
The last bound and the head bound follow by substituting
\(r\rho^m\le(\alpha\cdot16)^m/8800000\), \(n-5\ge91\), and the lower
bound on \(\pi\). The insertion costs at most \(4Js/n<s/2\).

\paragraph{Exterior \(x\ge3\sqrt p\).}
Use \(\mathcal H_j,T_j(h),S_j,W\) as in Section~\ref{sec:six}.
For \(n\ge96\),
\[
 rn^7<16,\qquad D_6(n)\le n^6/600,\qquad 6n^7<2^{n-7}.
\]
For the first inequality \(\sqrt{95}>487/50\), and the Taylor
sum through degree 96 at \(3(487/50)\) exceeds \(96^7/16\).
The function \(y^7e^{-3\sqrt{y-1}}\) decreases for \(y\ge96\)
since \(9y>14^2\). For the other two inequalities use the decrease of
\(D_6(n)/n^6\), and check the power-of-two bound at 96, then use
\((97/96)^7<2\).

For \(h\ge1\), \(j\le n\), coefficient extraction gives
\(|T_j(h)/h^j-1|\le e^{\mathcal V(n)}-1<1/(6n)\), since
\(\mathcal V(n)(6n+1)\le 16(6+1/n)/600<1\).
At \(h=0\) and \(n-7\le j\le n\) use the nonconstant tilt mass at \(2n\).
Here \(\mathcal V(2n)\le64rn^6/600<1/2\), and \(j!\le n^j\) gives
\[
 |T_j(0)|\le 2^{-j}\bigl(e^{\mathcal V(2n)}-1\bigr)
          \le 2^{-j}(128/600)rn^6<r/(6n).
\]
In the modal inequalities proved in Section~\ref{sec:six}, the
\(p\)-th absolute mass is at most \(4W\), the \(n\)-th at most
\(nW/2\), and
\(W/2<|S_n|+(m_p+1)|S_p|\) with \(m_p+1\le n/8\).
They apply to both indices because \(3\sqrt{n-1}\ge(29/10)\sqrt{n}\).
The two label bounds then give
\[
 |\mathcal H_j-S_j|<W/n\quad(n-7\le j\le p),\qquad
 |\mathcal H_n-S_n|<W/8 .
\]
At \(\mathcal H_p=0\) this implies \(|\mathcal H_n|>W/4\).
The shorter rows have \(|\mathcal H_j|<2W\) by unimodality of the
unsigned logistic weights. Their total contribution is at most
\(2W(21/20)rD_6(n)<W/4\). This proves the strict gap in the exterior.

\paragraph{The first 88 pairs for the independent hard-cutoff proof.}
The closed analytic intervals meet at \(x=16\) and \(x=3\sqrt p\).
Together they give \eqref{goal} throughout the five-dimensional cube
for every \(n\ge96\). The reciprocal and derangement arguments in
Section~\ref{transport} supply the original root order and simplicity.

For hard cutoff seven and \(8\le n\le95\), compute the two polynomial
rows independently by choosing the cycle through the smallest label
and by inserting the largest label. The first method uses the recurrence
\[
 E_{n,7}=\sum_{a=7}^{n}\binom{n-1}{a-1}tA_{a-1}(t)E_{n-a,7},
 \qquad E_{0,7}=1 .
\]
The second uses the insertion recurrence for the original rows from
\cite[Theorem~30]{rr}. Literal permutations check sizes 7, 8 and 9.
For each pair the certificate places all \(n-3\) lower roots in strictly
ordered negative dyadic intervals \([a_i,b_i]\), with opposite endpoint
signs. The upper row takes the same nonzero sign at \(a_i\) and \(b_i\),
alternating between consecutive intervals. It also has a certified
opposite sign to the left of the first interval, and its constant term
has sign opposite to the last bracket sign. Thus its \(n-2\) roots occur
strictly in the gaps and on both sides of the bracket list.
The exact degrees exhaust the roots of both rows and prove simplicity.
The data contain 4,268 lower-root brackets and 88 left-hand points.
This completes Theorem~\ref{seven}.
